\documentclass[12pt, b5paper]{article}
\setlength{\topmargin}{-.5in}
\setlength{\textwidth}{6.5in}
\setlength{\textheight}{9.5in}
\setlength{\oddsidemargin}{-.25in}

\usepackage{amsmath}

\newcommand{\floor}[1]{\left\lfloor #1 \right\rfloor}

\pagestyle{empty}
\begin{document}

\begin{center}
{\bf CSC 226
Exercise 3}\\
An Nguyen, V01073373\\
31 Jan 2026
\end{center}

\section*{Problem:}
Consider a hash table with size m (indices  $0, 1, ..., m-1$) which uses 
chaining to resolve collisions. Suppose the following weird hash function 
is used:
$$ h(k) =  1 + (k\text{ mod }(m-1)) $$
Notice that the function $h(k)$ does not satisfy the uniform hashing 
assumption, since no element will ever hash to index zero. Suppose n randomly 
selected keys are inserted into the table. Derive a formula for the expected 
number of items in the list at index 1 of the table and justify your derivation. 


\section*{Answer:}
From the question, the hash function $h(k)$ above does not assign any keys to index $0$,
so it does not satisfy the Uniform Hashing Assumption (UHA). However, as mentioned in lecture,
we can view the UHA as an assumption of the behaviour of an idealized hash function in our analysis, since the 
benefits of uniform hashing apply to some extent even in cases where a hash function is 
\textit{almost} uniform (which is ours).
\subsection*{Claim:} Under the Uniform Hashing Assumption, the expected size of the list $L_j$ at each 
index $j$ of the table is $\frac{n}{m}$.
\subsection*{Justification:} 
Consider a randomly selected sequence $k_1, \ldots, k_n$ of keys in $U$. Define a random variable $X_{ij}$ as follows:
$$
X_{ij} = 
\begin{cases}
    1 & \text{if } h(k_i) = j \\
    0 & \text{otherwise } 
\end{cases}
$$
By the UHA, we have: $$ P(X_{ij} = 1) = P(h(k_i) = j) = 1/m 
\text{ for } 1 \leq i \leq n \text{ and } 0 \leq j \leq m.
$$
Since we are trying to find the expected number if items in the list at index 1, and 
considering the hash function: 
$$ h(k) = 1 + (k \text{ mod } (m-1)), $$
we can imply that a key hashes to index 1 if and only if $k \text{ mod } (m-1) = 0$.\\
Values of $(k\text{ mod }(m-1))$ range from $0, 1, \ldots, m-2$, thus: 
$$ P(h(k_1)=1) = \frac{1}{m-1} $$
Therefore, the expected value of $X_{i1}$ is:
$$ E[X_{i1}] = 1 \cdot P(X_{i1} = 1) + 0 \cdot P(X_{i1} = 0) = \frac{1}{m-1} $$
The size of $L_1$ is then equal to the number of keys hashed to index $1$, which is
$$K_1 = X_{11} + X_{21} + \ldots + X_{n1}$$
so the expected size of $L_1$ is: 
$$E[K_1] = E[X_{11} + X_{21} + \ldots + X_{n1}] = E[X_{11}] + \ldots + E[X_{n1}] = n \cdot \frac{1}{m-1}$$
Hence, the desired formula for expected size of the list $L_j$ at index $1$: 
 \[
 \boxed{E[K_1] =\frac{n}{m-1}} 
 \]
\end{document}